% ════════════════════════════════════════════
% 模型的改进与推广（8.模型改进推广.tex）写作规则 —— 模板内嵌，AI 先读本文件再写
% ════════════════════════════════════════════
\section{模型的改进与推广}

\subsection{模型的改进}

本文模型可在方法、参数与数据三个层面进一步改进。

方法层面，可将配比—损失的线性混合回归替换为 Dirichlet 回归，或引入领域间交互项与稀疏先验（如 Group Lasso），以捕捉非线性效应并把 13 个域平均 $R^2=0.629$ 中未被解释的约 $37\%$ 变异进一步压缩；对未测量域（nih\_exporter 等四个域）可尝试用域文本统计特征构造代理变量，替代当前“固定于参考配比”的处理；冲突消解可由当前的“多数指标方向裁决”升级为基于潜变量模型（如有序 probit 因子模型）的一致性度量，从而给出冲突样本的质量可信度区间而非二值标记。

参数层面，广义标度律的质量缺口项可放宽为更灵活的形式（如 $C(1-Q)^{\gamma_1}D^{-\gamma_2}$ 或质量与数据量的交互项），以检验“质量缺口是否与数据规模耦合”；算力优化模型可进一步纳入数据重复轮次（epoch）、稀疏化/量化带来的推理与训练成本差异、以及分阶段课程学习等新决策变量；质量成本函数可在附录 B 三型之外引入由真实清洗流水线标定的经验成本曲线，以缓解低预算档对成本函数形式较敏感的问题。

数据层面，可用更多具备质量标注的公开语料扩充 B6/B7 的质量区间（当前外推下限受数据支持范围限制），并尝试从语料元数据（发布时间、来源许可、去重比例）反推质量标签；对附件 C4 可通过组织名、参数量与发布时间做模糊匹配，与 C1 建立更细粒度的行级连接，从而用真实训练算力直接约束前沿预测；对问题四，若能获得跨评测标准（v1 与 v2）的校准样本，则可把 2019--2023 的历史锚点统一到同一标尺，显著缓解当前“一致窗口仅 10 个月”的外推局限。

\subsection{模型的推广}

本文提出的“损失函数—成本约束—资源配置”建模框架具有广泛的推广价值。

在工程领域，该框架可直接迁移至算力受限场景下的大模型预训练、微调与蒸馏资源配置决策：式\eqref{eq:p3_balance}给出的最优性条件与“等效附加数据量 $\tau(Q)$”这一折算量，使得数据治理投入与算力扩张投入可以直接放在同一成本口径下比较，帮助企业以最小算力达到目标能力。

在数据治理领域，本文的质量综合评价与冲突消解模型（含 22 指标的显式方向表与“冲突率—独立基准”对比框架）可用于通用语料清洗、数据筛选与领域配比的动态调整；本文提出的可替代度指标 $s_{ij}$（式中的系数向量余弦相似度）可用于识别哪些领域可以合并、哪些领域必须搭配使用，为配比策略提供结构化依据。

在行业与政策层面，本文“长周期 shift-share 分解 + 近端一致窗口分解”的双重口径方法，可为科研机构评估“堆参数”与“改进算法”两条路线的投入产出、制定算力与数据资源的宏观配置政策提供定量依据；同时本文在方法上刻意区分“可直接使用的稳健结论”“存在临界点的结论”与“需谨慎表述的结论”，这一证据分级范式对涉及多源异构数据的政策评估同样适用。

更进一步，能力前沿的饱和外推方法（logistic 曲线外推与结构外推双路线互为校验）可推广至芯片制程、新能源技术、生物医药研发等具有渐进饱和特征的技术演进预测场景；而本文对半合成数据所做的“方向审计 + 不可约损失下限检验”自洽性筛查流程，可作为使用半合成/估算数据开展定量建模时的通用质量把关步骤。
